\section{Notation, Definitions, and General Assumptions for Proofs}
\label{app:notation-for-proofs}

We recall some definitions and define additional notation used in the proofs.
Additional terms are defined in the appendix sections in which they are used.

The heterogeneous treatment effects $\bm\tau$ are drawn from the prior distribution $\mathcal{P}$.

Conditional on the heterogeneous treatment effects $\bm\tau$, the stage-$t$ average treatment effect is
\[
  \ATE_t = \frac{\bm s_t^\top \bm\tau}{n_t},
  \qquad t\in\{1,2,3\},
\]
and the pilot and follow-up estimates satisfy
\[
  \widehat{\ATE}_t = \ATE_t+\frac{\varepsilon_t}{\sqrt{n_t}},
  \qquad
  \varepsilon_t\sim\mathcal N(0,1),
  \qquad
  t\in\{1,2\},
\]
with the stage-specific sampling noises mutually independent of each other and of $\bm\tau$.
We also write the conditional pass probability for the significance test at stage $t$ as
\[
  p_t(x)\defeq \Phi\!\big(\sqrt{n_t}\,x-z_{1-\alpha_t}\big)\in(0,1),
  \qquad t\in\{1,2\},
\]
so that $p_t(\ATE_t)=\P(\Test_t=1\mid \bm\tau)$ almost surely. Here $\Phi$ represents the normal cumulative distribution function. We will write $\varphi$ for the normal density, and $\phi$ for the generating characteristic function of a generic elliptical distribution.

Define the four first-stage objectives discussed in \Cref{sec:alternate-objectives-appendix} as
\begin{align*}
  V_{\singlesuccess}(\bm s_1) &\defeq \E[\Test_1],\\
  V_{\twosuccess}(\bm s_1) &\defeq \E[\Test_1\cdot\Test_2],\\
  V_{\singleimpact}(\bm s_1) &\defeq \E[\Test_1\cdot \ATE_3],\\
  V_{\twoimpact}(\bm s_1) &\defeq \E[\Test_1\cdot\Test_2\cdot \ATE_3].
\end{align*}
where throughout much of the paper body we write the main two-stage impact objective as
\[
  V(\bm s_1)\defeq V_{\twoimpact}(\bm s_1).
\]
For the two-stage impact objective define the oracle impact $V_{\oracle}$, the noise-free pilot impact $V_{\noisefree}$, and the unconditional follow-up impact $V_{\followup}$:
\begin{align*}
  V_{\oracle}
  &\defeq \E\!\left[\indic{\ATE_3>0}\cdot\Test_2\cdot \ATE_3\right],\\
  V_{\noisefree}(\bm s_1)
  &\defeq \E\!\left[\indic{\ATE_1>0}\cdot\Test_2\cdot \ATE_3\right],\\
  V_{\followup}
  &\defeq \E\!\left[\Test_2\cdot \ATE_3\right].
\end{align*}

For a random downstream payoff $Y$, define the corresponding first-stage objective
\[
  V_Y(\bm s_1)\defeq \E[\Test_1\cdot\, Y].
\]
We will say that $Y$ is \emph{non-adaptive} if it is independent of the pilot-stage noise conditional on $\bm\tau$:
\[
  Y \ind \varepsilon_1 \mid \bm\tau.
\]
For such $Y$, define
\[
  \theta_d^Y\defeq \E[\tau_d\cdot Y],
  \qquad
  \pi_d^Y\defeq \frac{\theta_d^Y}{\sqrt{c_d}}
\]
and write $\bm\theta^Y = (\theta_d^Y)_{d=1}^D$.

The four named objectives correspond to the downstream payoffs
\[
  Y_{\singlesuccess}\defeq 1,
  \qquad
  Y_{\twosuccess}\defeq \Test_2,
  \qquad
  Y_{\singleimpact}\defeq \ATE_3,
  \qquad
  Y_{\twoimpact}\defeq \Test_2\cdot \ATE_3,
\]
all of which are non-adaptive.
For the main two-stage impact objective we retain the shorthand
\[
  \theta_d \defeq \theta_d^{\twoimpact}
  = \E[\tau_d\cdot\Test_2\cdot \ATE_3],
  \qquad
  \idx_d = \pi_d^{\twoimpact} = \frac{\theta_d}{\sqrt{c_d}},
\]
and write $\bm\theta = (\theta_d)_{d=1}^D$.

Recall the definition of the set of unit-investment normalized pilot designs from \Cref{sec:Model},
\[
  \overline{\mathcal{S}} \defeq \{\bm x\in\R_{\ge0}^D:\bm c^\top\bm x=1\},
\]
and define the unit-investment representative pilot design as
\[
  \overline{\bm s}_1^{\,\infty}
  \defeq
  \frac{\bm s_3}{\bm c^\top \bm s_3}\in\overline{\mathcal{S}}.
\]
